\editionPage[\pdfbookmark[0]{Comprendre pourquoi}{bm-preuves-debut-parent}]{preuves-debut}{Comprendre pourquoi}{Pourquoi multiplier par 10 décale les chiffres}{Les explications ci-dessous éclairent des règles utilisées dans le livre.}
\begin{blocEdition}{preuves-debut}
\editionEtape{Revenir aux rangs}Dans 47, le 4 compte des dizaines et le 7 des unités. Multiplier par 10, c’est rendre chaque paquet dix fois plus grand : les 7 unités deviennent 7 dizaines, les 4 dizaines deviennent 4 centaines.\editionEtape{Lire le résultat}Le nombre obtenu a donc un 4 aux centaines, un 7 aux dizaines et plus rien aux unités : on écrit 470. Chaque chiffre a monté d’un rang, et le zéro garde la place des unités.\editionEtape{Et pour les décimaux}De même, $3{,}7\times10$ transforme 3 unités en 3 dizaines et 7 dixièmes en 7 unités : on obtient 37. La virgule semble glisser, mais ce sont les chiffres qui changent de rang.\editionRetenir{« Ajouter un zéro » n’est pas la règle : elle échoue sur $3{,}7$. La vraie règle est le changement de rang de chaque chiffre.}
\end{blocEdition}
\editionPage{preuves-2}{Comprendre pourquoi}{Pourquoi on peut ajouter des parts de même taille}{Les explications ci-dessous éclairent des règles utilisées dans le livre.}
\begin{blocEdition}{preuves-2}
\editionEtape{Une même unité}Sur une bande partagée en huit, chaque part est un huitième de la bande : toutes les parts ont exactement la même taille.\label{fig-preuves-2-preuve-fractions}\figureGeo{preuve-fractions}
\editionEtape{Réunir}Prendre 3 parts puis 2 parts, c’est prendre 5 parts en tout : $\frac38+\frac28=\frac58$. On compte des objets identiques, comme 3 billes et 2 billes font 5 billes.\editionEtape{Pourquoi le dénominateur ne change pas}Le 8 décrit le partage de la bande, qui n’a pas bougé : seul le nombre de parts prises a changé.\editionRetenir{Pour additionner des fractions, il faut des parts de même taille, donc le même dénominateur. Les numérateurs s’ajoutent, le partage reste.}
\end{blocEdition}
\editionPage{preuves-3}{Comprendre pourquoi}{Pourquoi le périmètre du carré vaut quatre fois le côté}{Les explications ci-dessous éclairent des règles utilisées dans le livre.}
\begin{blocEdition}{preuves-3}
\editionEtape{Suivre le contour}Faire le tour d’un carré, c’est parcourir ses quatre côtés l’un après l’autre : le périmètre est la somme côté + côté + côté + côté.\editionEtape{Regrouper}Quatre termes identiques s’écrivent avec une multiplication : côté $\times4$. Pour un côté de 7 cm, le tour vaut $4\times7=28$ cm.\editionEtape{Adapter à d’autres figures}Le même raisonnement donne la règle du triangle équilatéral (3 fois le côté) ou du rectangle (deux longueurs et deux largeurs). La règle n’est pas à apprendre par cœur : elle se retrouve en suivant le contour.\editionRetenir{Comprendre d’où vient une règle permet de la reconstruire au besoin, et d’éviter de la confondre avec une formule d’aire.}
\end{blocEdition}
\editionPage{preuves-4}{Comprendre pourquoi}{Pourquoi deux nombres pairs ont une somme paire}{Les explications ci-dessous éclairent des règles utilisées dans le livre.}
\begin{blocEdition}{preuves-4}
\editionEtape{Dire ce qu’est un nombre pair}Un nombre pair se range exactement en paquets de deux, sans objet isolé : 26, c’est 13 paquets de deux.\editionEtape{Réunir deux collections}Additionner 26 et 58, c’est réunir 13 paquets et 29 paquets : on obtient 42 paquets de deux, sans rien d’isolé. La somme 84 est donc paire.\editionEtape{Un argument pour tous les cas}Ce raisonnement ne dépend pas des nombres choisis : réunir des paquets de deux donne toujours des paquets de deux. La propriété est vraie pour tous les nombres pairs, pas seulement pour les exemples testés.\editionRetenir{Un exemple vérifie ; un raisonnement sur « tous les cas » explique et prouve. Un seul contre-exemple suffirait, au contraire, à démolir une règle annoncée.}
\end{blocEdition}
